\chapter{CORS, sessioni, filtri e autenticazione}

\section{Cross-Origin Resource Sharing (CORS)}

\begin{definizione}[Origine e CORS]
L'\textbf{origine} di un'applicazione client-side \`e il server HTTP che l'ha fornita.
Si parla di \textbf{cross-origin resource sharing} quando un'app client-side vuole effettuare
una fetch da una app server-side che risiede su un'\emph{origine diversa}.
\end{definizione}

\subsection{La preflight request}
Quando si verifica una richiesta cross-origin, il browser -- \emph{prima} di far partire la
fetch vera e propria -- effettua una richiesta detta \textbf{preflight}, con method HTTP
\textbf{OPTIONS}: un test preliminare per capire se la richiesta effettiva potr\`a essere
recepita dal server. Nella preflight il browser descrive la richiesta che vorrebbe fare con
tre header:
\begin{itemize}
  \item \texttt{Access-Control-Request-Method}: quale sar\`a il method;
  \item \texttt{Origin}: l'origine dell'app client-side;
  \item \texttt{Access-Control-Request-Headers}: gli header della richiesta che verr\`a
        effettuata.
\end{itemize}
Il server risponde alla preflight dichiarando i propri vincoli con gli header:
\begin{itemize}
  \item \texttt{Access-Control-Allow-Methods}: i method permessi;
  \item \texttt{Access-Control-Allow-Headers}: gli header permessi;
  \item \texttt{Access-Control-Allow-Origin}: le origini permesse;
  \item \texttt{Access-Control-Allow-Credentials}: se \`e permesso l'invio di credenziali
        (essenzialmente i cookie di sessione).
\end{itemize}
Il browser cos\`i sa se la richiesta verr\`a accettata e quali header dovr\`a inserire.
\textbf{Se la preflight viene respinta, \`e il browser stesso a bloccare la richiesta
originale}; se la mandasse avanti, sarebbe l'app server-side a rispondere con un errore
(es. \texttt{405 Method Not Allowed} o \texttt{403 Forbidden} a seconda della violazione).

\subsection{@CrossOrigin}
Per default i server HTTP \textbf{bloccano} le richieste cross-origin (politica
\emph{noCORS}). \`E un problema in fase di sviluppo, quando abbiamo \emph{due} server HTTP
(uno serve l'app client-side, l'altro \`e embedded nell'app server-side): anche sullo stesso
host (localhost) girano su \emph{porte diverse} $\Rightarrow$ origini diverse.

\textbf{Approccio semplificato} (da usare \emph{solo in sviluppo su localhost}): annotare con
\textbf{@CrossOrigin} i controller (o i singoli handler) per cui permettere le richieste
cross-origin -- in sviluppo, tipicamente tutti. Effetto: quando la richiesta \`e diretta a
uno dei nostri controller viene permessa \emph{indipendentemente} da origine, method e
header -- una sorta di ``salvacondotto universale''. Solo l'invio di \textbf{credenziali}
resta bloccato. (Le richieste che non corrispondono a nessun handler restano soggette a
noCORS: vengono bloccate in preflight prima ancora di stabilire se sono valide.)

\section{Le sessioni}

\begin{definizione}[Sessione]
Quando c'\`e un'interazione prolungata fra due ``parti'' (utente e applicazione, client e
server, applicazioni peer-to-peer) diventa rilevante il concetto di \textbf{sessione}:
l'idea che la ``conversazione'' abbia \emph{uno stato e uno storico}, cos\`i che a ogni passo
si possa fare riferimento a quanto gi\`a detto e fatto, senza ripetere o ricapitolare.
\end{definizione}

\begin{importante}[Perch\'e serve uno stratagemma]
Il problema nasce da due ragioni:
\begin{enumerate}
  \item il protocollo HTTP \`e di per s\'e \textbf{stateless}: non prevede alcuna idea di
        stato dello scambio -- ogni transazione (richiesta + risposta) \`e indipendente
        dalle altre; l'unico modo di ``simulare'' uno stato sarebbe riportare a ogni passo
        tutte le informazioni rilevanti gi\`a scambiate;
  \item un server HTTP gestisce \textbf{molte conversazioni con client diversi}: se il client
        non si rende riconoscibile, il server non pu\`o sapere ``a che punto'' della
        conversazione si trova, perch\'e \`e in punti diversi a seconda dell'interlocutore.
\end{enumerate}
Serve quindi ``costruire'' \emph{sopra} HTTP un meccanismo che permetta a client e server di
sapere che stanno avendo una conversazione prolungata e di mantenerne lo stato.
\end{importante}

\subsection{Sessioni tramite cookie}
\begin{definizione}[Cookie]
Un \textbf{cookie} \`e un blocco di dati che il server manda al client \emph{nel primo
scambio} della loro conversazione; a ogni scambio successivo, il client \emph{aggiunge il
cookie alla propria richiesta}, cos\`i il server identifica il particolare ``thread'' di
conversazione con quel client.
\end{definizione}
Il server \emph{potrebbe} memorizzare le informazioni direttamente nel cookie; di solito
per\`o nel cookie si mettono solo le informazioni per \emph{riconoscere l'utente e
identificare univocamente la conversazione} (un id di sessione), mentre \textbf{lo stato vero
e proprio resta sul server} (in memoria o nel DB), indicizzato da quell'id: cos\`i si evita
di appesantire richieste e risposte HTTP.

\subsection{Gestione della sessione in Spring Boot}
Per attivare i cookie di sessione si configura \texttt{application.properties}:
\begin{lstlisting}
server.servlet.session.cookie.name=LAB_SESSION_COOKIE
server.servlet.session.cookie.http-only=true
server.servlet.session.cookie.secure=true
\end{lstlisting}
\begin{itemize}
  \item il \emph{nome} identifica il cookie (visibile negli strumenti di sviluppo del
        browser);
  \item \textbf{http-only}: il browser e il codice client-side \emph{non hanno accesso} ai
        contenuti del cookie;
  \item \textbf{secure}: il cookie viene trasmesso \emph{cifrato}.
\end{itemize}
Per lavorare con la sessione si \emph{inietta nei request handler} l'oggetto
\textbf{HttpSession}, tramite cui:
\begin{itemize}
  \item scrivere: \texttt{session.setAttribute(key, value)};
  \item leggere: \texttt{session.getAttribute(key)};
  \item invalidare (anche prima del timeout): \texttt{session.invalidate()};
  \item impostare il timeout: \texttt{session.setMaxInactiveInterval(secs)} -- il timeout
        indica per quanto tempo il cookie pu\`o restare \emph{inutilizzato} (nessuna nuova
        richiesta con quel cookie); scaduto, il cookie non \`e pi\`u valido e dovr\`a
        essere generato uno nuovo.
\end{itemize}

\subsection{La collaborazione del client (e il CORS con credenziali)}
Perch\'e il meccanismo funzioni, l'app client-side deve \emph{rimandare il cookie a ogni
richiesta}. \`E il comportamento di default del browser \textbf{solo per richieste
same-origin}: in situazione \emph{cross-origin}, per default le credenziali \emph{non}
vengono inviate. Per ottenerne l'invio, nelle opzioni della fetch si specifica:
\begin{lstlisting}[language=JavaScript]
credentials: "include"      // default: "same-origin"
\end{lstlisting}
Lato server, per\`o, \texttt{@CrossOrigin} cos\`i com'\`e non basta pi\`u:
\begin{enumerate}
  \item la risposta alla preflight deve includere
        \texttt{Access-Control-Allow-Credentials: true} $\Rightarrow$
        \texttt{@CrossOrigin(allowCredentials = "true")};
  \item per ragioni di sicurezza, \textbf{se le credenziali sono ammesse il protocollo HTTP
        impedisce di abilitare tutte le origini in modo acritico}: l'impostazione di default
        (\texttt{Allow-Origin: qualsiasi}) non \`e pi\`u valida -- vanno \emph{elencate
        esplicitamente} le origini ammesse (una \textbf{whitelist}):
\end{enumerate}
\begin{lstlisting}[style=java]
@CrossOrigin(origins = "http://localhost:5173", allowCredentials = "true")
\end{lstlisting}

\section{I filtri}

\begin{definizione}[Filtro]
Un \textbf{filtro} in Spring Boot/Jakarta \`e una funzione attraverso cui passano le
richieste HTTP \emph{prima} di giungere ai request handler: viene chiamata per \emph{ogni}
richiesta, e pu\`o decidere se \emph{inoltrarla} nella \textbf{filter chain} (la catena di
filtri, che culmina nell'invocazione dell'handler) oppure \emph{respingerla} con uno status
code di errore. Uso tipico: verificare che la richiesta soddisfi dei requisiti.
\end{definizione}

Un filtro \`e una classe che \textbf{implementa l'interfaccia \texttt{Filter}} ed \`e
annotata \textbf{@Component}: cos\`i Spring Boot lo include nella filter chain (che
comprende anche filtri gi\`a presenti nel framework). L'ordine di esecuzione si controlla
con \textbf{@Order(N)}.
\begin{lstlisting}[style=java]
@Component
@Order(1)
public class AuthFilter implements Filter {
  @Override
  public void doFilter(ServletRequest servletRequest,
                       ServletResponse servletResponse,
                       FilterChain filterChain)
                       throws IOException, ServletException {
    HttpServletRequest request = (HttpServletRequest) servletRequest;
    HttpServletResponse res = (HttpServletResponse) servletResponse;
    ...
    filterChain.doFilter(servletRequest, servletResponse); // inoltro
    ...
    res.sendError(HttpServletResponse.SC_FORBIDDEN);       // blocco
  }
}
\end{lstlisting}
I parametri di \texttt{doFilter}:
\begin{itemize}
  \item \texttt{ServletRequest}: la richiesta nella sua forma ``grezza'', prima che Spring la
        elabori; rappresenta una generica richiesta di rete -- per le informazioni specifiche
        HTTP serve il \emph{type casting} a \texttt{HttpServletRequest};
  \item \texttt{ServletResponse}: la risposta ``grezza'', in cui Spring inserir\`a le
        informazioni degli handler; anche qui casting a \texttt{HttpServletResponse};
  \item \texttt{FilterChain}: la catena a cui inoltrare i due oggetti, per passare il
        testimone al prossimo filtro o, se i filtri sono finiti, all'handler.
\end{itemize}
Per \textbf{inoltrare}: \texttt{filterChain.doFilter(request, response)}. Per
\textbf{bloccare}: non si invoca la chain e si usa la response per inviare un errore:
\texttt{response.sendError(HttpServletResponse.SC\_XXX)}.

\section{Basic Authentication con filtri e sessioni}

Con filtri e sessioni si pu\`o ``cucinare'' un meccanismo di \textbf{autenticazione di
base}. Versione semplificata (e per nulla sicura): username e password viaggiano \emph{in
chiaro come query parameters}; una versione un po' pi\`u sicura li mette nel \emph{body} di
una POST e usa \textbf{HTTPS}.

Schema:
\begin{itemize}
  \item si distingue un utente autenticato da uno non autenticato tramite un
        \textbf{attributo salvato in sessione} (es. lo username);
  \item per una generica richiesta, \textbf{il filtro verifica se la sessione contiene
        l'attributo}: se s\`i, la fa passare; altrimenti risponde
        \texttt{401 UNAUTHORIZED};
  \item \textbf{eccezione: la richiesta di login}, che deve passare in ogni caso. Arriva al
        request handler, che (tramite un opportuno @Service) verifica le credenziali: se
        corrette, crea in sessione l'attributo che segnala l'autenticazione; altrimenti
        risponde \texttt{401 UNAUTHORIZED};
  \item il \textbf{logout} si realizza semplicemente \emph{invalidando la sessione};
  \item il server deve esporre anche una route, accessibile pure ai non autenticati, che
        permetta all'app client-side di \emph{verificare lo stato corrente
        dell'autenticazione}.
\end{itemize}
